\section{Weighted trees and a canonical response field}
\label{sec:response}

The analytic input is a response estimate whose constant is independent
of the number of edges and of their weights.  We first prove it on finite
graphs, then use the free spanning-tree limit to construct a
drift on every field of real edge marks. The finite determinant and
transfer-current identities are classical; see Burton and
Pemantle~\cite{BurtonPemantle1993} and the weighted formulation in
\cite[Section~4]{BLPS2001}. We prove the needed identities and their
free-limit consequence below.

\begin{lemma}[Weighted-tree response]
\label{lem:tree-response}
Let $H$ be a finite connected graph with at least two vertices, and let
$\mathcal T(H)$ denote its set of spanning trees.  For
$h\in\mathbb R^{E(H)}$, give each $T\in\mathcal T(H)$ probability
proportional to $\exp(\sum_{f\in T}h_f)$, and write
$\xi_e=\1_{\{e\in T\}}$ and $p_e(h)=\mathbb E_h\xi_e$.
For every $e\in E(H)$,
\begin{equation}
\label{eq:covariance-row}
\frac{\partial p_e}{\partial h_f}
 =\operatorname{Cov}_h(\xi_e,\xi_f),
\qquad
\sum_{f\in E(H)}
 \left|\frac{\partial p_e}{\partial h_f}\right|
 =2p_e(h)(1-p_e(h))\le\frac12.
\end{equation}
In particular, for all $h,h'\in\mathbb R^{E(H)}$,
\begin{equation}
\label{eq:tree-response}
 |p_e(h)-p_e(h')|
 \le\frac12\max_{f\in E(H)}|h_f-h'_f|.
\end{equation}
\end{lemma}

\begin{proof}
Set $c_f=e^{h_f}$.  Orient the edges temporarily, delete one row of the
oriented incidence matrix, and denote the resulting columns by $a_f$.
The reduced weighted Laplacian is
\[
 L=\sum_{f\in E(H)}c_fa_fa_f^{\mathsf T}.
\]
Write $r=|V(H)|-1$.  Cauchy--Binet, applied to the reduced incidence
matrix with column $f$ multiplied by $\sqrt{c_f}$, gives
\[
 \det L
 =\sum_{\substack{F\subseteq E(H)\\|F|=r}}
   \det(a_f:f\in F)^2\prod_{f\in F}c_f
 =\sum_{T\in\mathcal T(H)}\prod_{f\in T}c_f
 =:\mathcal Z_H(h).
\]
For the second equality, a set containing a cycle has dependent
incidence columns, while an acyclic set of $r$ edges is a spanning
tree.  Its reduced incidence determinant has absolute value one:
expand at a leaf other than the vertex whose row was deleted and
continue by induction.  Connectedness supplies a spanning tree, so
$\mathcal Z_H(h)>0$ and $L$ is invertible.

Differentiating the finite partition sum gives both
$p_e=\partial_{h_e}\log\mathcal Z_H$ and the covariance identity in
\eqref{eq:covariance-row}.  Differentiating the determinant instead gives
\[
 p_e=c_ea_e^{\mathsf T}L^{-1}a_e.
\]
Since
$\partial_{h_f}L^{-1}=-c_fL^{-1}a_fa_f^{\mathsf T}L^{-1}$,
for distinct edges $e,f$ we obtain
\begin{equation}
\label{eq:negative-square}
 \operatorname{Cov}_h(\xi_e,\xi_f)
 =-c_ec_f\bigl(a_e^{\mathsf T}L^{-1}a_f\bigr)^2\le0.
\end{equation}
Every spanning tree has exactly $r$ edges.  Consequently
\[
 \sum_{f\in E(H)}\operatorname{Cov}_h(\xi_e,\xi_f)
 =\operatorname{Cov}_h\!\left(\xi_e,
                  \sum_{f\in E(H)}\xi_f\right)=0.
\]
The diagonal covariance is $p_e(1-p_e)$; all the other covariances
are nonpositive by \eqref{eq:negative-square}.  Their absolute sum
is therefore $2p_e(1-p_e)\le1/2$, as claimed.  This also covers a
bridge: its inclusion probability is one and its entire covariance
row is zero.  Finally, integrate the gradient along the segment from
$h$ to $h'$ to obtain \eqref{eq:tree-response}.  The estimate holds
at every point of that segment, without a restriction on the weights.
\end{proof}

\subsection{The free limit from finite determinants}

We include the free-limit construction, so that the passage from finite
weighted trees to the infinite law has an explicit foundation. The relevant
space is the closure of the finite cycle flows; no description of all
square-summable divergence-free flows on the infinite graph is needed.
This cycle-space description is classical; see
\cite[Proposition~7.1 and Theorem~7.8]{BLPS2001}.

\begin{proposition}[Free exhaustion limit]\label{prop:free-limit}
Let $G$ be an infinite connected locally finite simple graph. If
$H_1\subseteq H_2\subseteq\cdots$ are finite connected subgraphs with
$\bigcup_n E(H_n)=E(G)$, the uniform spanning-tree laws on $H_n$, extended
by absent edges outside $H_n$, converge weakly on $\{0,1\}^{E(G)}$.
Their limit is independent of the exhaustion, is supported on spanning
forests, and is transported by every graph isomorphism. We denote this
law by $\FUSF_G$.
\end{proposition}

\begin{proof}
Temporarily orient all edges of $G$. For a finite connected subgraph $H$
with at least two vertices, let $A_H$ be its reduced incidence matrix and
put
\[
 K_H=A_H^{\mathsf T}(A_HA_H^{\mathsf T})^{-1}A_H.
\]
For each finite $F\subseteq E(H)$, the matrix-tree calculation in
Lemma~\ref{lem:tree-response}, followed by the rectangular determinant
identity, gives the polynomial identity
\begin{align}
 \E\prod_{e\in F}(1+z_e\xi_e)
 &=\frac{\det(A_HA_H^{\mathsf T}
        +A_F\operatorname{diag}(z)A_F^{\mathsf T})}
        {\det(A_HA_H^{\mathsf T})} \notag\\
 &=\det\!\left(I_F+\operatorname{diag}(z)(K_H)_{F,F}\right).
 \label{eq:finite-inclusion-polynomial}
\end{align}
Here $A_F$ consists of the columns indexed by $F$. The first equality
can first be read for $z_e>-1$ as a change of edge weights, and then
holds for all $z$ because both sides are polynomials. The identity
$\det(I+UV)=\det(I+VU)$ used in the second line follows by eliminating
blocks in $\left(\begin{smallmatrix}I&U\\-V&I\end{smallmatrix}\right)$
in either order. Expanding the last determinant in its diagonal variables
and comparing coefficients yields
\begin{equation}\label{eq:finite-inclusion-minor}
 \P(F\subseteq T_H)=\det (K_H)_{F,F}.
\end{equation}
This includes sets larger than $|V(H)|-1$: the corresponding minor and
tree-inclusion probability are both zero. No minor is inverted.

Regard finite edge vectors as vectors of $\ell^2(E(G))$ by zero extension.
Let $C_H=\ker A_H$ in the coordinates of $E(H)$, and let $Q_H$ be the
orthogonal projection onto those coordinates. The matrix $K_H$ is the
orthogonal projection onto the row space of $A_H$, so its zero extension is
\begin{equation}\label{eq:finite-cycle-projection}
 K_H=Q_H-P_{C_H}.
\end{equation}
Deleting a different incidence row does not change $C_H$: the full
incidence matrix has entries summing to zero in each column, so zero
divergence at all but one vertex
forces zero divergence at the remaining vertex. The space $C_H$ is exactly
the span of its oriented cycle flows. Indeed, choose a spanning tree of
$H$ and subtract the fundamental cycle flow for each non-tree edge to
eliminate that edge's coefficient. The remaining circulation is supported
on a tree and vanishes by successive leaf removal.

Write $C_n=C_{H_n}$. These spaces increase under zero extension: a
circulation supported on $H_n$ still has zero divergence at every vertex
of $H_{n+1}$. Every finite cycle of $G$ is eventually contained in $H_n$,
because the exhaustion includes every edge. Therefore the closure
\[
 C=\overline{\bigcup_n C_n}
\]
is the closed span of all finite cycle flows of $G$, independent of the
chosen exhaustion, whether or not its members are induced subgraphs.
The projections $P_{C_n}$ converge strongly to $P_C$. To see this directly,
for $m\ge n$ and $x\in\ell^2(E(G))$ orthogonality gives
\[
 \|P_{C_m}x-P_{C_n}x\|^2
 =\|P_{C_m}x\|^2-\|P_{C_n}x\|^2.
\]
The squared norms increase and are bounded by $\|x\|^2$, so the projected
vectors are Cauchy. Their limit lies in $C$, and its difference from $x$
is orthogonal to every $C_n$, hence to $C$. This identifies the limit as
$P_Cx$. Also $Q_{H_n}\to I$ strongly. Thus the zero-extended $K_{H_n}$
converge strongly to $K=I-P_C$, and every fixed finite principal minor
converges to that of $K$.

For disjoint finite edge sets $A,B$, inclusion--exclusion now gives a limit
for
\[
 \P(A\subseteq T_{H_n},\ B\cap T_{H_n}=\varnothing)
 =\sum_{J\subseteq B}(-1)^{|J|}
       \det(K_{H_n})_{A\cup J,A\cup J}
\]
once $A\cup B\subseteq E(H_n)$; the empty determinant is one.
These limits are nonnegative, normalized and consistent on every finite
cube, since they are limits of actual probabilities. To construct their
countable law explicitly, temporarily enumerate the edges and recursively
split $[0,1)$ into half-open intervals of the prescribed binary-prefix
masses. A uniform point selects one successive prefix at every depth and
hence a configuration with these finite marginals. Zero-mass intervals
may be omitted. Cylinder probabilities uniquely determine the resulting
law, so the temporary enumeration affects neither the law nor the graph
rule constructed later. Continuous functions on the compact countable
binary product are uniformly approximable by finite-coordinate functions;
cylinder convergence consequently proves weak convergence.

For every finite cycle, the event that all its edges are present has
probability zero in every sufficiently late tree law and hence in the
limit. There are only countably many finite cycles, so the limit is
acyclic on the full vertex set. Reversing any orientations conjugates each
finite principal matrix by a diagonal matrix with entries $\pm1$ and
preserves its determinant. Graph isomorphisms transport the finite cycle
space and all the limiting cylinder probabilities. This proves the claimed
independence and invariance.
\end{proof}

\subsection{Finite tilts and the canonical drift}

Fix now an infinite connected locally finite simple undirected graph
$G$.  Its edge set $E=E(G)$ is countable.  Write
$\nu=\FUSF_G$ for the free uniform spanning forest law, viewed as a law
on $\{0,1\}^{E}$.  For $e\in E$, finite
$S\subseteq E$, $t\ge0$, and $y\in\mathbb R^{E}$, define
\begin{equation}
\label{eq:finite-tilt}
 b_{e,S}(t,y)
 =\frac{\displaystyle
    \int x_e\exp\!\left(\sum_{f\in S}(y_f-t/2)x_f\right)\,d\nu(x)}
   {\displaystyle
    \int \exp\!\left(\sum_{f\in S}(y_f-t/2)x_f\right)\,d\nu(x)}.
\end{equation}
The denominator is positive, since its integrand is at least
$\exp(\sum_{f\in S}\min\{y_f-t/2,0\})$.

\begin{lemma}[Finite-field forest response]
\label{lem:forest-response}
The function $b_{e,S}$ takes values in $[0,1]$ and is continuous in
$(t,(y_f)_{f\in S})$.  For all $t\ge0$ and $y,y'\in\mathbb R^{E}$,
\begin{equation}
\label{eq:finite-response}
 |b_{e,S}(t,y)-b_{e,S}(t,y')|
 \le\frac12\max_{f\in S}|y_f-y'_f|,
\end{equation}
where the maximum over the empty set is defined to be zero.
\end{lemma}

\begin{proof}
Both integrals in \eqref{eq:finite-tilt} are finite sums over the
possible configurations on $S\cup\{e\}$.  Positivity of the
denominator proves continuity, and the range follows from $0\le x_e\le1$.
If $S$ is empty, the function is the constant $\nu(x_e=1)$.

For nonempty $S$, take a finite connected free exhaustion $H_m$ of $G$.
For all sufficiently large $m$, it contains $S\cup\{e\}$.  Give its
edges log weights $y_f-t/2$ for $f\in S$ and zero elsewhere.
The resulting tree inclusion probability is the ratio in
\eqref{eq:finite-tilt} with the uniform spanning-tree law on $H_m$
in place of $\nu$.  Lemma~\ref{lem:tree-response} bounds the difference
between this probability and the one using $y'$ by the right side of
\eqref{eq:finite-response}.  With $S,t,y$ fixed, both integrands
are bounded cylinder functions, so their expectations converge under
the free spanning-tree limit.  The limiting denominator is
strictly positive by the displayed lower bound.  Thus the ratios
converge to $b_{e,S}(t,y)$, and likewise for $y'$.  Passing to the
limit proves \eqref{eq:finite-response}.
\end{proof}

To specify one drift at every field, we now fix the outer exhaustion
intrinsically.  If $e=\{u,v\}$, let $H_n(e)$ be the induced subgraph
on
\[
 \{z\in V(G):\min\{d_G(z,u),d_G(z,v)\}\le n\},
 \qquad n\ge1,
\]
and set $S_n(e)=E(H_n(e))$.  These graphs are finite and connected,
increase with $n$, and exhaust $G$.  Define, for every $t\ge0$ and
every $y\in\mathbb R^{E}$,
\begin{equation}
\label{eq:drift}
 b_e(t,y)=\limsup_{n\to\infty}b_{e,S_n(e)}(t,y),
 \qquad b(t,y)=(b_e(t,y))_{e\in E}.
\end{equation}

\begin{proposition}[Canonical drift]
\label{prop:canonical-drift}
Each $b_e$ is a jointly Borel function on
$[0,\infty)\times\mathbb R^{E}$ with values in $[0,1]$, where
$\mathbb R^{E}$ has its product Borel structure.  The construction
commutes with graph isomorphisms and does not select a root.  For every
$t\ge0$ and $y,y'\in\mathbb R^{E}$ such that
$\sup_{f\in E}|y_f-y'_f|<\infty$, it satisfies
\begin{equation}
\label{eq:response-bound}
 \sup_{e\in E}|b_e(t,y)-b_e(t,y')|
 \le\frac12\sup_{f\in E}|y_f-y'_f|.
\end{equation}
Neither $y$ nor $y'$ is required to have bounded coordinates.
\end{proposition}

\begin{proof}
The range and Borel assertions follow from
Lemma~\ref{lem:forest-response} and the countable limsup in
\eqref{eq:drift}.  If $d=\sup_f|y_f-y'_f|<\infty$, that lemma gives
\[
 b_{e,S_n(e)}(t,y)
 \le b_{e,S_n(e)}(t,y')+d/2
\]
for every $e,n$.  Taking limsups and then interchanging $y,y'$ proves
\eqref{eq:response-bound}.

For the intrinsic nature of the construction, the inner limit has the
explicit expression
\begin{equation}
\label{eq:canonical-finite-formula}
 b_{e,S_n(e)}(t,y)
 =\lim_{m\to\infty,\ m\ge n}
 \frac{\displaystyle
  \sum_{T\in\mathcal T(H_m(e))}
   \1_{\{e\in T\}}
   \exp\!\left(\sum_{f\in T\cap S_n(e)}(y_f-t/2)\right)}
 {\displaystyle
  \sum_{T\in\mathcal T(H_m(e))}
   \exp\!\left(\sum_{f\in T\cap S_n(e)}(y_f-t/2)\right)}.
\end{equation}
This is exactly the fixed-finite-field limit used in the proof of
Lemma~\ref{lem:forest-response}.  Each finite ratio is invariant under
isomorphisms preserving the distinguished undirected edge and its
field marks.  The prescribed limits preserve that property.  The same
finite formulas therefore specify the drift as the graph varies;
their use in a Borel graph rule will be made explicit below.
\end{proof}

The order of the two spatial limits matters: in
\eqref{eq:canonical-finite-formula}, $n$ is fixed while $m$ tends to
infinity, and only afterwards is the limsup in $n$ taken.  We neither
differentiate an infinite-volume partition function nor use a
cardinality identity for an infinite forest.  We also do not require
the outer limit to exist at arbitrary fields.  The limsup supplies a
bounded Borel version everywhere, and
\eqref{eq:response-bound} controls its change under every bounded
perturbation of a possibly unbounded field.
